\NativeTitle{101}{期末综合素养评价（三）}{}{90}{5}{一 & 二 & 三 & 四 & 五}{ &  &  &  & }
\NativeSlot{279}{40}{%
\ExamGroup{一、填空。{\normalfont\kaishu\fontsize{10.8bp}{14bp}\selectfont （每空1分，共15分）}}

% p101_s01

\ExamQuestion{1.}{5}%
一个数是由\(6\)个一、\(3\)个十分之一和\(5\)个百分之一组成，这个数写作\mbox{（\hspace{4.4em}）}，读作\mbox{（\hspace{7em}）}，精确到十分位是\mbox{（\hspace{1.8em}）}。\par
}
\NativeSlot{354}{36}{%
% p101_s02

\ExamQuestion{2.}{5}%
把\(128+147=275\)，\allowbreak{}\(275\div25=11\)，\allowbreak{}\(240\times11=2640\)这三个算式写成一个综合算式是\mbox{（\hspace{13em}）}。\par
}
\NativeSlot{406}{50}{%
% p101_s03

\ExamQuestion{3.}{5}%
在○里填上“＞”“＜”或“＝”。\par \(0.060\) ○ \(0.06\)；\allowbreak{}\(7\)角\(5\)分 ○ \(0.8\)元；\allowbreak{}\(7.2\)米 ○ \(7\)米\(2\)厘米；\allowbreak{}\(65\times100-1\) ○ \(65\times99\)；\allowbreak{}\(1000-156-244\) ○ \(1000-(156+244)\)\par
}
\NativeSlot{493}{36}{%
% p101_s04

\ExamQuestion{4.}{5}%
一个等腰三角形，它的一个底角是\(45^\circ\)，它的顶角是\mbox{（\hspace{1.8em}）}°，按角分类，这个三角形是\mbox{（\hspace{1.8em}）}三角形。\par
}
\NativeSlot{543}{36}{%
% p101_s05

\ExamQuestion{5.}{5}%
把一张正方形纸沿相同的方向对折两次，打开后出现两条折痕，折痕的位置关系是\mbox{（\hspace{6em}）}。\par
}
\NativeSlot{594}{57}{%
% p101_s06

\ExamQuestion{6.}{5}%
如图，一张三角形和一张平行四边形纸片。把三角形的各个角撕下来拼在一起，可以组成一个\mbox{（\hspace{1.8em}）}角，把平行四边形的各个角撕下来拼在一起可以组成一个\mbox{（\hspace{1.8em}）}角。\par
\BookFigure{figures/p101_s06-01.tex}
}
\NativeSlot{697}{36}{%
% p101_s07

\ExamQuestion{7.}{5}%
四年级一班参加书法社团的有\(18\)人，参加绘画社团的有\(16\)人，两个社团都参加的有\(5\)人，只参加一个社团的有\mbox{（\hspace{1.8em}）}人。\par
}
